% LSPV-3 REPORT::PREDTYPE1::2026-10-05T06:00:00-07:00
\documentclass{article}
\usepackage{amsmath}
\begin{document}
\section*{PREDTYPE1: acquisition versus return typing}
\textbf{STATUS: OPEN.} P versus NP remains open.
\section*{CORPUS REVIEWED}
Reviewed governance, CODEX, pnp-research skill, root main, MAIN RETURNTRACE1, ACQ1, RETURNSQUEEZE1, RECURCLOSE1, NORMBLOW1, SHORTKERN1, DUALSHORT1-R2, SUCCINCTESCAPE1, PHABSORB1, FULLSKOLEM1, active tasks, and HEC/CRS-PPC/CCSPC. Repository search checked prior acquisition and return formulations.
\section*{PROVED/EXACT type separation}
Boundary-Squeeze requires an exact repeatable semantic edge from an original SAT instance F to a returned SAT instance Q with
\[F\in SAT\iff Q\in SAT.\]
ACQ1 instead defines the NP predicate
\[P_q(v)=L(v)\wedge[B_vq\ne0],\]
whose compilation asks whether a legal state violates q. No checked source theorem equates that auxiliary predicate with satisfiability of the original F. Therefore a short SAT compilation of P_q is an acquisition query, not a RECURCLOSE1 return query. Its exponent cannot be substituted for the return-size exponent without a separate satisfiability-preserving theorem.
\section*{Charge separation}
For degree two, M=1+m+\binom m2. Dense q has order m squared written bits. A conditional O(m)-support relation can have a sparse descriptor near O(m log m), but sparse-Q normalization remains OPEN. Short acquisition and semantic return are distinct requirements. If acquisition makes n^u queries of length n^beta to a SAT decider of exponent K, charge u+beta K plus deterministic costs. Only a separately proved return edge may contribute eta+theta K to recursive closure.
\section*{Quantifier and literature audit}
The full reconstruction predicate remains fixed-level alternating PH syntax. Conditional PH absorption under P=NP is qualitative and does not provide a fractional return compiler. Cook--Levin preserves the NP predicate actually compiled; it does not make an auxiliary NP predicate equivalent to a different SAT instance without a reduction theorem. SAT time-space bounds measure actual encoded SAT length.
\section*{DUPLICATION / CONTRADICTION AUDIT}
ACQ1 already proves NP typing and charged acquisition. RETURNTRACE1 already records that no ordinary SAT return object is source-certified. This report adds the explicit type guard; it does not contradict either result.
\section*{WHAT CHANGED RELATIVE TO ROOT MAIN}
Root main already requires an exact return relation. The new correction explicitly excludes ACQ1 as a substitute absent a semantic theorem.
\section*{ADOPTED / REJECTED / RETIRED PRIOR COMPONENTS}
Adopt ACQ1 for acquisition and RECURCLOSE1 only conditionally after a true return edge exists. Retire the shortcut of computing a return zeta directly from P_q. DEAD/COPY/DELAY/INTERP/RECOMP and semantic/algebraic oracles do not create the missing equivalence; this channel remains relativizing.
\section*{STALE OR OVERSTATED CLAIM FOUND}
The previous proposal to compute verifier exponent tau for P_q, set zeta=b tau, and test zeta<1 as a Boundary-Squeeze return criterion is overstated. It is acquisition accounting only until original-SAT satisfiability preservation is proved.
\section*{Next target}
Search the exact vMS reconstruction proof for a source-specific NP-typed predicate equivalent to the original SAT instance, with no current-scale hidden SAT call. If none exists, prove a source-stop lemma before further exponent optimization.
\section*{Status}
PROVED/EXACT: predicate-type separation. FALSIFIED/RETIRED: ACQ1 exponent as return exponent without a semantic theorem. OPEN: exact SAT-semantic return, sparse-Q normalization, arbitrary-algorithm transfer, nonrelativization, nonalgebrization, and P versus NP.
% END_OF_VERIFIED_REPORT::LSPV-3::PREDTYPE1
\end{document}
